\documentclass[12pt]{article}
\usepackage[T1]{fontenc}
\usepackage[utf8]{inputenc}
\usepackage{amsmath,amssymb,array,booktabs,longtable,geometry,enumitem}
\geometry{letterpaper,margin=0.90in}
\setlength{\parindent}{1.25em}
\setlength{\parskip}{0.34em}
\setlength{\emergencystretch}{3em}
\allowdisplaybreaks
\sloppy
\begin{document}

\begin{center}
{\Large C. F. Gauss}\\[3pt]
{\large \emph{Theory of Biquadratic Residues. Second Memoir}}\\[3pt]
{\large Title page and Articles 24--29, printed pp. 93--101}
\end{center}

\section*{Theory of Biquadratic Residues. Second Memoir\\
Title page and Articles 24--29}

\begin{center}
\textsc{Theory of Biquadratic Residues}\\
\textsc{Second Memoir}\\[2pt]
\textsc{By Carl Friedrich Gauss}\\
\textsc{Presented to the Royal Society, 15 April 1831}\\
\emph{Commentationes Societatis Regiae Scientiarum Gottingensis recentiores}, Vol. VII.\\
Göttingen, 1832.
\end{center}

\subsection*{24.}

In the first memoir, all that is required for the biquadratic classification of the number $+2$ was completely settled. We divide all numbers not divisible by a modulus $p$, where $p$ is a prime of the form $4n+1$, among four complexes $A,B,C,D$ according as they become congruent modulo $p$ to $+1,+f,-1,-f$ when raised to the exponent $\frac14(p-1)$, where $f$ denotes either root of $ff\equiv-1\pmod p$. We found that the decision as to the complex to which $+2$ belongs depends on the decomposition of $p$ into two squares. Thus, if $p=a^2+b^2$, where $a^2$ is the odd square and $b^2$ the even square, and if the signs of $a$ and $b$ are chosen so that $a\equiv1\pmod4$ and $b\equiv af\pmod p$, then $+2$ belongs to $A,B,C,D$ according as $\frac12 b$ has the form $4n,4n+1,4n+2,4n+3$, respectively.

The rule for classifying $-2$ follows at once. Since $-1$ belongs to $A$ when $\frac12 b$ is even, and to $C$ when it is odd, Article 7 gives that $-2$ belongs to $A,B,C,D$ according as $\frac12 b$ has the form $4n,4n+3,4n+2,4n+1$, respectively. Equivalently:
\[
\begin{array}{c|cc}
\text{Complex} & +2 & -2\\
\hline
A & b\equiv 0 & b\equiv 0\\
B & b\equiv 2a & b\equiv 6a\\
C & b\equiv 4a & b\equiv 4a\\
D & b\equiv 6a & b\equiv 2a
\end{array}
\qquad(\bmod\,8).
\]
It is easy to see that, stated in this way, the theorems no longer depend on the condition $a\equiv1\pmod4$; they remain valid also when $a\equiv3\pmod4$, provided the other condition $af\equiv b\pmod p$ is retained. The whole result can also be compressed into the single formula
\[
b^{\frac12ab}\equiv a^{\frac12ab}\,2^{\frac14(p-1)}\pmod p,
\]
when $a$ and $b$ are taken positive.

\subsection*{25.}

We next ask how far induction indicates the classification of the number $3$. Continuing the table of Article 11, always taking the least primitive root, shows that $+3$ belongs to
\[
\begin{array}{c|rr@{\quad}c|rr@{\quad}c|rr@{\quad}c|rr}
\multicolumn{3}{c}{A\ \text{for}}&
\multicolumn{3}{c}{B\ \text{for}}&
\multicolumn{3}{c}{C\ \text{for}}&
\multicolumn{3}{c}{D\ \text{for}}\\
p&a&b&p&a&b&p&a&b&p&a&b\\\hline
13&-3&+2&17&+1&-4&37&+1&-6&5&+1&+2\\
109&-3&+10&29&+5&+2&61&+5&-6&41&+5&-4\\
181&+9&+10&53&-7&+2&73&-3&-8&149&-7&+10\\
193&-7&-12&89&+5&-8&97&+9&+4&173&+13&+2\\
229&-15&+2&101&+1&+10&157&-11&-6&&&\\
277&+9&+14&113&-7&-8&241&-15&-4&&&\\
&&&137&-11&-4&&&&&&\\
&&&197&+1&-14&&&&&&\\
&&&233&+13&+8&&&&&&\\
&&&257&+1&-16&&&&&&\\
&&&269&+13&+10&&&&&&\\
&&&281&+5&+16&&&&&&\\
&&&293&+17&+2&&&&&&
\end{array}
\]
At first sight no simple connection appears among the values $a,b$ that correspond to the same complex. But if we recall that, in the theory of quadratic residues, the analogous decision is simpler for $-3$ than for $+3$, there is hope for an equally successful outcome here. We find that $-3$ belongs to
\[
\begin{array}{c|rr@{\quad}c|rr@{\quad}c|rr@{\quad}c|rr}
\multicolumn{3}{c}{A\ \text{for}}&
\multicolumn{3}{c}{B\ \text{for}}&
\multicolumn{3}{c}{C\ \text{for}}&
\multicolumn{3}{c}{D\ \text{for}}\\
p&a&b&p&a&b&p&a&b&p&a&b\\\hline
37&+1&-6&5&+1&+2&13&-3&+2&29&+5&+2\\
61&+5&-6&17&+1&-4&73&-3&-8&41&+5&-4\\
157&-11&-6&89&+5&-8&97&+9&+4&53&-7&+2\\
193&-7&-12&113&-7&-8&109&-3&+10&101&+1&+10\\
&&&137&-11&-4&181&+9&+10&197&+1&-14\\
&&&149&-7&+10&229&-15&+2&269&+13&+10\\
&&&173&+13&+2&241&-15&-4&293&+17&+2\\
&&&233&+13&+8&277&+9&+14&&&\\
&&&257&+1&-16&&&&&&\\
&&&281&+5&+16&&&&&&
\end{array}
\]
and here the inductive law presents itself immediately:
\[
\begin{array}{l|l}
A& b\equiv0\pmod3\\
B& a+b\equiv0\pmod3,\ \text{that is } b\equiv2a\pmod3\\
C& a\equiv0\pmod3\\
D& a-b\equiv0\pmod3,\ \text{that is } b\equiv a\pmod3.
\end{array}
\]

\subsection*{26.}

The number $+5$ is found to belong to
\[
\begin{array}{l|l}
A& p=101,109,149,181,269\\
B& p=13,17,73,97,157,193,197,233,277,293\\
C& p=29,41,61,89,229,241,281\\
D& p=37,53,113,137,173,257.
\end{array}
\]
Taking into account the values of $a,b$ belonging to each $p$, the law is just as easy to grasp as for the classification of $-3$:
\[
\begin{array}{l|l}
A& b\equiv0\pmod5\\
B& b\equiv a\pmod5\\
C& a\equiv0\pmod5\\
D& b\equiv4a\pmod5.
\end{array}
\]
These rules plainly cover all cases; for if $b\equiv2a$ or $b\equiv3a\pmod5$, then $a^2+b^2\equiv0$, contrary to the hypothesis that $p$ is a prime different from $5$.

\subsection*{27.}

Applying induction in the same way to the numbers $-7,-11,+13,+17,-19,-23$, and carrying it far enough, gives the following rules:
\[
\begin{array}{c|l}
\multicolumn{2}{c}{\text{For }-7}\\\hline
A& a\equiv0,\ \text{or }b\equiv0\pmod7\\
B& b\equiv4a,\ \text{or }5a\\
C& b\equiv a,\ \text{or }6a\\
D& b\equiv2a,\ \text{or }3a
\end{array}
\]
\[
\begin{array}{c|l}
\multicolumn{2}{c}{\text{For }-11}\\\hline
A& b\equiv0,5a,\ \text{or }6a\pmod{11}\\
B& b\equiv a,3a,\ \text{or }4a\\
C& a\equiv0,\ \text{or }b\equiv2a,\ \text{or }9a\\
D& b\equiv7a,8a,\ \text{or }10a
\end{array}
\]
\[
\begin{array}{c|l}
\multicolumn{2}{c}{\text{For }+13}\\\hline
A& b\equiv0,4a,9a\pmod{13}\\
B& b\equiv6a,11a,12a\\
C& a\equiv0;\ b\equiv3a,10a\\
D& b\equiv a,2a,7a
\end{array}
\]
\[
\begin{array}{c|l}
\multicolumn{2}{c}{\text{For }+17}\\\hline
A& a\equiv0;\ b\equiv0,a,16a\pmod{17}\\
B& b\equiv2a,6a,8a,14a\\
C& b\equiv5a,7a,10a,12a\\
D& b\equiv3a,9a,11a,15a
\end{array}
\]
\[
\begin{array}{c|l}
\multicolumn{2}{c}{\text{For }-19}\\\hline
A& b\equiv0,2a,5a,14a,17a\pmod{19}\\
B& b\equiv3a,7a,11a,13a,18a\\
C& a\equiv0;\ b\equiv4a,9a,10a,15a\\
D& b\equiv a,6a,8a,12a,16a
\end{array}
\]
\[
\begin{array}{c|l}
\multicolumn{2}{c}{\text{For }-23}\\\hline
A& a\equiv0;\ b\equiv0,7a,10a,13a,16a\pmod{23}\\
B& b\equiv2a,3a,4a,11a,15a,17a\\
C& b\equiv a,5a,9a,14a,18a,22a\\
D& b\equiv6a,8a,12a,19a,20a,21a.
\end{array}
\]

\subsection*{28.}

The special theorems obtained by induction in this way are confirmed as far as the induction is continued, and they exhibit a very beautiful form of criteria. When they are compared in order to draw general conclusions from them, the following observations occur at once.

The criteria for deciding to which complex the prime number $\pm q$ is to be referred, taking the upper or lower sign according as $q$ has the form $4n+1$ or $4n+3$, depend on the forms of $a,b$ compared with one another modulo $q$.

First, when $a\equiv0\pmod q$, $\pm q$ belongs to a fixed complex: namely to $A$ for $q=7,17,23$, and to $C$ for $q=3,11,13,19$. This suggests that the former case holds in general when $q$ is of the form $8n\pm1$, and the latter when $q$ is of the form $8n\pm3$. Moreover, $B$ and $D$ are excluded without induction when $a$ is divisible by $q$, since then $p\equiv b^2\pmod q$, so that $p$ is a quadratic residue of $q$, and therefore by the fundamental theorem $\pm q$ must be a quadratic residue of $p$.

Second, when $a$ is not divisible by $q$, the criterion depends on the value of $\frac{b}{a}\pmod q$. This expression admits $q$ different values, $0,1,2,\ldots,q-1$; but when $q$ is of the form $4n+1$, the two values of $\sqrt{-1}\pmod q$ must be excluded, since they plainly cannot be values of $\frac{b}{a}$ under the standing assumption that $p=a^2+b^2$ is a prime distinct from $q$. Thus the number of admissible values of $\frac{b}{a}\pmod q$ is $q-2$ for $q\equiv1\pmod4$, while it remains $q$ for $q\equiv3\pmod4$.

These values are distributed into four classes. Let certain values, denoted indefinitely by $\alpha$, correspond to $A$; others, denoted by $\beta$, to $B$; others, $\gamma$, to $C$; and the remaining values, $\delta$, to $D$. Thus $\pm q$ is assigned to $A,B,C,D$ according as $b\equiv\alpha a$, $b\equiv\beta a$, $b\equiv\gamma a$, $b\equiv\delta a\pmod q$. The law of this distribution appears more hidden, although certain general facts can be noticed at once. Three of the classes have the same number of values, $\frac14(q-1)$ or $\frac14(q+1)$, while in one class, namely the class corresponding to the complex with criterion $a\equiv0$, the number is smaller by one. Hence the total number of distinct criteria for each complex is the same, namely $\frac14(q-1)$ or $\frac14(q+1)$.

We also observe that $0$ is always in the first class, among the $\alpha$, and that the complements to $q$ of the numbers $\alpha,\beta,\gamma,\delta$ lie respectively in the first, fourth, third, and second classes. Finally, the values of $1/\alpha,1/\beta,1/\gamma,1/\delta\pmod q$ are seen to belong to the first, fourth, third, and second classes when the criterion $a\equiv0$ corresponds to $A$; but to the third, second, first, and fourth classes, respectively, when that criterion is referred to $C$. This is about the limit of what induction can provide, unless we boldly anticipate results that will below be drawn from their genuine sources.

\subsection*{29.}

Before proceeding further, it is useful to observe that the criteria for prime numbers, taken positive when they are of the form $4n+1$ and negative when they are of the form $4n+3$, are enough to decide the case of all remaining numbers, provided Article 7 and the criteria for $-1$ and $\pm2$ are also used. For example, if criteria are wanted for $+3$, the criteria given in Article 25 for $-3$ remain valid for $+3$ when $\frac12b$ is even; but when $\frac12b$ is odd, the complexes $A,B,C,D$ must be permuted with $C,D,A,B$. Hence:
\[
\begin{array}{c|l}
\multicolumn{2}{c}{+3\ \text{belongs}}\\\hline
A& b\equiv0\pmod{12};\ \text{or also }a\equiv0\pmod3,\ b\equiv2\pmod4\\
B& b\equiv8a,\ \text{or }10a\pmod{12}\\
C& b\equiv6a\pmod{12};\ \text{or also }a\equiv0\pmod3,\ b\equiv0\pmod4\\
D& b\equiv2a,\ \text{or }4a\pmod{12}.
\end{array}
\]
Likewise the criteria for $\pm6$ are obtained by combining the criteria for $\mp2$ and $-3$:
\[
\begin{array}{c|l}
\multicolumn{2}{c}{+6\ \text{belongs}}\\\hline
A& b\equiv0,2a,22a\pmod{24};\ \text{or also }a\equiv0\pmod3,\ b\equiv4a\pmod8\\
B& b\equiv4a,6a,8a\pmod{24};\ \text{or also }a\equiv0\pmod3,\ b\equiv2a\pmod8\\
C& b\equiv10a,12a,14a\pmod{24};\ \text{or also }a\equiv0\pmod3,\ b\equiv0\pmod8\\
D& b\equiv16a,18a,20a\pmod{24};\ \text{or also }a\equiv0\pmod3,\ b\equiv6a\pmod8,
\end{array}
\]
whereas $-6$ belongs as follows:
\[
\begin{array}{c|l}
\multicolumn{2}{c}{-6\ \text{belongs}}\\\hline
A& b\equiv0,10a,14a\pmod{24};\ \text{or also }a\equiv0\pmod3,\ b\equiv4a\pmod8\\
B& b\equiv4a,8a,18a\pmod{24};\ \text{or also }a\equiv0\pmod3,\ b\equiv6a\pmod8\\
C& b\equiv2a,12a,22a\pmod{24};\ \text{or also }a\equiv0\pmod3,\ b\equiv0\pmod8\\
D& b\equiv6a,16a,20a\pmod{24};\ \text{or also }a\equiv0\pmod3,\ b\equiv2a\pmod8.
\end{array}
\]
In the same way, criteria for $+21$ will be assembled from the criteria for $-3$ and $-7$; criteria for $-105$ from those for $-1,-3,+5,-7$, and so on.

\end{document}
